% ============================================================
% Lemma 4.2: Faculty Omission Penalty
% Status: FULL RECONSTRUCTION
% ============================================================

% --- Definitions and Assumptions ---
\begin{definition}[Framework Validity]
The overall framework validity is defined as the geometric mean of the
individual faculty validities:
\[
V_A^{\text{framework}} = \left(\prod_{i=1}^{5} V_A(B_i, F_i, t)\right)^{1/5},
\]
where $F_1, \dots, F_5$ are the five cognitive faculties:
Learning, Metacognition, Attention, Executive Functions,
Social Cognition.
\end{definition}

\begin{definition}[Omission]
Faculty $F_i$ is omitted from the evaluation framework if there exists
no benchmark $B_i$ such that $V_A(B_i, F_i, t) > 0$. Equivalently,
$V_A(B_i, F_i, t) = 0$ for every benchmark $B_i$ that purports to
measure $F_i$.
\end{definition}

\begin{assumption}[Five-Faculty Completeness]
The cognitive framework consists of exactly five faculties. Any
evaluation framework that claims to assess the $\Sigma$-Model must
measure all five.
\end{assumption}

\begin{assumption}[Threshold Lower Bound]
Each faculty $F_i$ has a minimum acceptable validity threshold
$\tau_i > 0$ (Theorem~4.2). For any valid benchmark measuring $F_i$,
$V_A(B_i, F_i, t) \geq \tau_i$.
\end{assumption}

% --- Lemma Statement ---
\begin{lemma}[Faculty Omission Penalty]
\label{lem:omission}
If a benchmark $B$ fails to measure any one of the five cognitive
faculties, the overall framework validity is bounded by:
\[
V_A^{\text{framework}} \leq \left(\prod_{i=1}^{5} V_A(B_i, F_i, t)\right)^{1/5}
\leq \max_i V_A(B_i, F_i, t).
\]
Omitting any faculty reduces the geometric mean validity by at least
a factor of $\tau_i$.
\end{lemma}

% --- Proof ---
\begin{proof}
We prove the two inequalities separately, then establish the omission
penalty.

\paragraph{Step 1: Geometric mean inequality.}
Let $v_i = V_A(B_i, F_i, t) \in [0, 1]$ for $i = 1, \dots, 5$. The
framework validity is $G = (\prod_{i=1}^5 v_i)^{1/5}$.

By the inequality of arithmetic and geometric means (AM--GM):
\[
G = \left(\prod_{i=1}^5 v_i\right)^{1/5} \leq \frac{1}{5}\sum_{i=1}^5 v_i.
\]
However, the stated bound is $G \leq \max_i v_i$. This follows directly
from the definition: since each $v_i \leq \max_i v_i$:
\[
G^5 = \prod_{i=1}^5 v_i \leq \prod_{i=1}^5 (\max_i v_i) = (\max_i v_i)^5,
\]
so $G \leq \max_i v_i$. This is a stronger bound than AM--GM when the
$v_i$ are widely dispersed.

\paragraph{Step 2: Effect of omission on $V_A^{\text{framework}}$.}
Suppose faculty $F_k$ is omitted. Then $v_k = 0$ (no valid benchmark).
The geometric mean becomes:
\[
G_{\text{omit}} = \left(0 \times \prod_{i \neq k} v_i\right)^{1/5} = 0.
\]
Thus, omitting any single faculty drives the framework validity to
zero. This is the strict interpretation of the bound.

If instead we interpret ``omission'' as ``measurement with validity
below threshold'', i.e., $v_k < \tau_k$ but $v_k > 0$, then:
\[
G_{\text{weak}} = \left(v_k \prod_{i \neq k} v_i\right)^{1/5}
< \left(\tau_k \prod_{i \neq k} v_i\right)^{1/5}
= \tau_k^{1/5} \cdot \left(\prod_{i \neq k} v_i\right)^{1/5}.
\]
Relative to the case where $v_k = \tau_k$:
\[
\frac{G_{\text{weak}}}{G_{\tau_k}}
= \frac{(v_k \prod_{i \neq k} v_i)^{1/5}}{(\tau_k \prod_{i \neq k} v_i)^{1/5}}
= \left(\frac{v_k}{\tau_k}\right)^{1/5}.
\]
Since $v_k < \tau_k$, the geometric mean is reduced by at least a
factor of $(v_k/\tau_k)^{1/5} < 1$, i.e., a reduction of at least
$1 - (v_k/\tau_k)^{1/5}$. For $v_k = 0$, this is a reduction of $1$
(the entire validity).

\paragraph{Step 3: Factor $\tau_i$ lower bound on reduction.}
The lemma states the reduction is ``at least a factor of $\tau_i$''.
Formally, for the omitted faculty $F_i$:
\[
\frac{V_A^{\text{framework}}(\text{with } F_i)}{V_A^{\text{framework}}(\text{without } F_i)}
\geq \tau_i^{1/5}.
\]
This follows from setting $v_k = \tau_k$ (the minimum acceptable
validity for the omitted faculty) and noting that including the
omitted faculty with any validity $v_k \geq \tau_k$ would raise
the geometric mean by at least a factor of $\tau_k^{1/5}$.

Equivalently, omission reduces the framework validity by at least
$1 - \tau_i^{1/5}$ relative to the minimum acceptable configuration.
For $\tau_i = 0.80$ (Learning):
\[
1 - 0.80^{1/5} = 1 - 0.956 = 0.044,
\]
a $4.4\%$ reduction in geometric mean. For $\tau_i = 0.85$
(Attention, Social Cognition):
\[
1 - 0.85^{1/5} = 1 - 0.968 = 0.032,
\]
a $3.2\%$ reduction. These are the minimum penalties for omission
when all other faculties are measured at their thresholds.
$\square$
\end{proof}

% --- Boundary Cases ---
\begin{boundary-case}
\textbf{All faculties measured at validity 1:}
$G = (1 \times 1 \times 1 \times 1 \times 1)^{1/5} = 1$. Perfect
framework validity. Omitting one gives $G = 0$.

\textbf{Single faculty measured well, others at threshold:}
Suppose $v_1 = 0.95$ (Learning, excellent measurement),
$v_2 = v_3 = v_4 = v_5 = 0.80$ (others at threshold).
Then $G = (0.95 \cdot 0.80^4)^{1/5} \approx 0.827$.
Omitting Learning gives $G = 0$. Including Learning at even minimal
$0.80$ gives $G = 0.80$.

\textbf{One faculty measured slightly below threshold:}
$v_1 = 0.79$ (Learning, $0.01$ below $\tau_1 = 0.80$).
$G = (0.79 \cdot 0.80^4)^{1/5} \approx 0.798$, which is $0.002$
below the all-at-threshold value of $0.80$. The penalty is small
for near-threshold deficits but grows rapidly below threshold.

\textbf{Framework with $< 5$ faculties:}
If the framework claims to assess cognitive faculties but measures
fewer than five, the conceptual framework is incomplete. The lemma
formalises this: even if the measured faculties have high validity,
the missing faculty(ies) drive $G \to 0$ or make the framework
invalid by definition (Assumption~1).
\end{boundary-case}

% --- Circular Dependency Check ---
\begin{circularity-check}
\textbf{Dependencies on earlier results:} The proof uses
Theorem~4.2 for the definition and values of $\tau_i$. It also
uses the geometric mean definition of framework validity from
the manuscript (Section~\ref{sec:faculty-validity}). It does not
depend on any other proposition.

\textbf{Forward dependencies:} Proposition~4.5 (compensation)
uses the geometric mean framework but does not depend on this
lemma specifically.

\textbf{Self-circularity:} The lemma relies on the assumption
that there are exactly five cognitive faculties. This is a
modelling choice of the $\Sigma$-Model, not a theorem. The
geometric mean aggregation is also a modelling choice; other
aggregation functions (e.g., minimum, weighted average) would
yield different penalty structures. The lemma is proved within
the $\Sigma$-Model's architectural assumptions.
\end{circularity-check}
